% Generated draft: requires mathematical and source review.
\documentclass[11pt]{article}
\usepackage[margin=25mm]{geometry}
\usepackage{fontspec}
\setmainfont{Latin Modern Roman}
\usepackage{amsmath,amssymb,mathtools,unicode-math}
\setmathfont{Latin Modern Math}
\usepackage{enumitem,needspace,xurl,hyperref}
\setlength{\parindent}{0pt}
\setlength{\parskip}{5pt}
\setlength{\emergencystretch}{3em}
\newcommand{\N}{\mathbb{N}}
\newcommand{\Z}{\mathbb{Z}}
\newcommand{\Q}{\mathbb{Q}}
\newcommand{\R}{\mathbb{R}}
\newcommand{\C}{\mathbb{C}}
\newcommand{\F}{\mathbb{F}}
\newcommand{\A}{\mathbb{A}}
\newcommand{\PP}{\mathbb{P}}
\newcommand{\E}{\mathbb{E}}
\newcommand{\eps}{\varepsilon}
\newcommand{\dd}{\,\mathrm{d}}
\begin{document}
\section*{UnsolvedMath ID: 30002334}
\textbf{Source problem code:} catalogue entry
\subsection{Definitions and mathematical statement}
For three distinct noncollinear points \(x,y,z\in\R^3\), let \(R(x,y,z)\) be the radius of their circumcircle and define their Menger curvature by
\[
 c(x,y,z)=\frac{1}{R(x,y,z)}
 =\frac{4\operatorname{Area}(\triangle xyz)}{|x-y|\,|y-z|\,|z-x|}.
\]
Set \(c(x,y,z)=0\) for collinear triples and for triples with repeated points. For a rectifiable geometric curve \(\Gamma\subset\R^3\) and \(p>0\), a conventional unrooted integral Menger energy is
\[
 \mathcal M_p(\Gamma)=\int_{\Gamma}\int_{\Gamma}\int_{\Gamma}
 c(x,y,z)^p\,\dd\mathcal H^1(x)\,\dd\mathcal H^1(y)\,\dd\mathcal H^1(z),
\]
where \(\mathcal H^1\) denotes one-dimensional Hausdorff measure. This definition supplies a conventional interpretation of the source notation \(M_p\); the supplied record does not specify its normalization, whether integration counts parametrization multiplicity, or its exact admissible curve class.

Fix \(L>0\), and let \(\mathcal A_L\) denote the intended admissible class of space curves of length \(L\). This class is unresolved in the supplied record: in particular, closedness, embeddedness, initial regularity, and any knot-class restriction are unstated. Subject to specifying this class and the energy convention, the source asks the following three questions.

\begin{enumerate}
\item Let \(C_L\subset\R^3\) be a round circle of circumference \(L\). For which intended exponents \(p\) does
\[
 \mathcal M_p(\Gamma)\geq \mathcal M_p(C_L)
 \qquad\text{for every }\Gamma\in\mathcal A_L
\]
hold, with equality only when \(\Gamma\) is a Euclidean isometric image of \(C_L\)? This is a template for the circle-minimization question, conditional on \(C_L\in\mathcal A_L\). The source does not state the range of \(p\), and uniqueness for parametrized curves would additionally require an equivalence under reparametrization.
\item At the critical exponent \(p=3\), does
\[
 M_3(\gamma)<\infty
 \quad\Longrightarrow\quad
 \gamma'(s)\text{ exists at every parameter value }s
\]
hold in the intended admissible class? To express this as an intrinsic tangent question, one would need to specify an arclength parametrization. For a closed curve this could have domain \(\R/L\Z\); for an open curve it could have domain \([0,L]\), with an explicit endpoint derivative convention. The supplied record specifies none of these choices, so the displayed implication is the source question with its unresolved parametrization and boundary qualifications, not a claim for arbitrary parametrizations.
\item What regularity must local minimizers and stationary points of \(M_p\) have? A precise formulation requires the intended exponents, admissible curves, topology defining local minimality, and permitted variations. It also requires specifying whether stationarity is constrained by the fixed length and whether additional boundary or knot constraints apply. No particular differentiability, H\"older, Sobolev, or smoothness conclusion is supplied as the target.
\end{enumerate}

These are statement-only questions. Their unresolved choices cannot be filled in uniquely from the supplied record.
\subsection{Short English statement}
The source asks whether a round circle uniquely minimizes integral Menger curvature at fixed length, whether finite critical energy \(M_3\) forces a tangent to exist everywhere, and how regular local minimizers and stationary curves must be. The admissible curves, exponent range, parametrization, boundary conditions, and variational framework require clarification before these become fully specified mathematical assertions.
\subsection{Sources}
\begin{itemize}
\item Local UnsolvedMath registry, supplied complete source record: permanent ID \(30002334\), title \emph{Critical Integral Menger Curvature}, exact source problem code \texttt{catalogue entry}, in the Oberwolfach Reports Open Problems collection. The record gives proposed year 2013 and supplies no named proposer.
\item Supplied report citation: \emph{Geometric Knot Theory (2014)}, Oberwolfach Reports, pp.\ 1313--1358, DOI \texttt{10.4171/owr/2013/22}; \url{https://doi.org/10.4171/owr/2013/22}. The background attributes the questions to workshop open-problem discussions. The same link is listed among the literature sources as \emph{Oberwolfach Report 22/2013}. The differing year designations are retained without resolution.
\item Strzelecki--von der Mosel, \emph{On some knot energies involving Menger curvature}; supplied literature link: \url{https://doi.org/10.1016/j.topol.2013.05.022}.
\item Strzelecki--Szumanska--von der Mosel, \emph{Regularizing and self-avoidance effects}; supplied literature link: \url{https://www.numdam.org/item/ASNSP_2010_5_9_1_145_0.pdf}.
\item Supplied literature assessment dated exactly \texttt{2026-08-21}: the record labels the problem \texttt{partially\_solved} and reports supercritical regularity and several variational results as known, while describing the critical \(M_3\) questions, circle uniqueness, and full minimizer or stationary-point regularity as open. This is supplied dated curation derived from the OWR material, not independently verified here; the record itself requests independent verification of time-sensitive status and cited claims.
\item The registry labels the statement \texttt{corrected\_verified} and reports a manual semantic, typographic, and LaTeX review of its cleanup. This is supplied editorial provenance and does not establish that the missing mathematical hypotheses have been resolved.
\end{itemize}
\end{document}
